% ============================================================================
% CHAPITRE 7 — CONCLUSION GÉNÉRALE ET PERSPECTIVES
% ============================================================================

\section*{Bilan du projet}
\addcontentsline{toc}{section}{Bilan du projet}

Le présent rapport de Projet de Fin d'Études rend compte de la conception, l'implémentation
et l'évaluation de \textbf{Guidni}, une plateforme de planification de voyage intelligente
dédiée à l'île de Djerba. En suivant rigoureusement la méthodologie Scrum sur quatre Sprints
(octobre 2024 -- février 2025), j'ai pu construire un système articulé autour de deux
sous-systèmes complémentaires.

Le premier sous-système, l'\textbf{Agent Planificateur conversationnel}, repose sur une
machine à états LangGraph à quatre nœuds (THINK, EXECUTE\_TOOL, VALIDATE, RESPOND) disposant
de dix-huit outils appelables par le LLM. Un pipeline RAG multilingue utilisant les embeddings
BAAI/bge-m3 (français, anglais, arabe) et un cross-encoder de réordonnancement garantissent
des réponses ancrées dans les données réelles du catalogue touristique de Djerba. L'abstraction
multi-fournisseurs LLM permet de basculer entre Ollama, Groq, Gemini et Lightning AI selon
les besoins. Le système de mémoire conversationnelle avec résumé automatique à 15~messages
assure la cohérence des conversations longues.

Le second sous-système, le \textbf{Moteur de Recommandation Hybride LinUCB}, exploite un
bandit contextuel avec un vecteur de caractéristiques de dimension~306 construit par produit
de Kronecker ($\phi = x_{\text{user}}(9) \otimes x_{\text{tags}}(34)$). Un pipeline de
classement post-traitement en dix étapes garantit la diversité, l'exploration et l'équité
des recommandations. Le cache matriciel en RAM avec verrou threading assure une latence de
scoring inférieure à 100~ms. Le tracker JavaScript vanilla collecte huit signaux
comportementaux en temps réel, alimentant l'apprentissage en ligne du modèle.

L'organisation Scrum s'est avérée déterminante pour tempérer les incertitudes inhérentes à
un projet d'intelligence artificielle. Les Sprint Reviews régulières ont permis de valider
chaque composant incrémentalement, et la Definition of Done a garanti un critère de qualité
objectif à chaque itération.


\section*{Perspectives et travaux futurs}
\addcontentsline{toc}{section}{Perspectives et travaux futurs}

Malgré les contributions apportées, plusieurs perspectives d'évolution se dégagent~:

\begin{enumerate}
    \item \textbf{Affinage d'un LLM local sur les données du tourisme tunisien.}
    L'historique de conversations stocké dans \texttt{PlannerMessage} constitue un corpus
    d'entraînement naturel. Un fine-tuning supervisé produirait un modèle domaine-spécifique
    comprenant la géographie locale et les niveaux de prix réalistes de Djerba.

    \item \textbf{Remplacement de LinUCB linéaire par un bandit neuronal.}
    Le modèle linéaire $r = \theta^\top \phi$ ne capture pas les interactions non linéaires
    complexes. Des architectures comme Neural UCB modéliseraient ces effets par un réseau de
    neurones tout en conservant la propriété UCB.

    \item \textbf{Infrastructure d'expérimentation A/B.}
    L'architecture actuelle ne permet pas de mesurer l'impact causal des changements
    algorithmiques. Un framework A/B permettrait de valider empiriquement les hypothèses
    d'amélioration avant déploiement.

    \item \textbf{Extension multi-régions.}
    Le champ \texttt{locationZone} est conçu pour supporter plusieurs zones géographiques.
    L'extension à Hammamet, Sousse et Tunis ne nécessite que l'ajout de données et
    l'initialisation des bras correspondants.

    \item \textbf{Interface vocale et multimodale.}
    L'intégration de la reconnaissance vocale arabe permettrait aux utilisateurs tunisiens
    de planifier leur voyage sans saisie clavier. Le pipeline RAG pourrait accepter des
    photos comme requêtes visuelles via des embeddings multimodaux.
\end{enumerate}


\section*{Mot de la fin}
\addcontentsline{toc}{section}{Mot de la fin}

Ce projet démontre que la combinaison de l'IA conversationnelle (agents LangGraph) et de
l'apprentissage en ligne (bandits contextuels) ouvre une nouvelle génération de plateformes
touristiques intelligentes~: simultanément personnalisées, adaptatives et ancrées dans des
données réelles. Le système Guidni est une preuve de concept opérationnelle démontrant que
cette combinaison est non seulement techniquement réalisable, mais pratiquement déployable
sur une infrastructure modeste~: un serveur unique avec Ollama local, une base de données
PostgreSQL standard et un tracker JavaScript sans aucune dépendance externe.
